% ED26-BM-RAT2021-10-4; Raub, MIT2021 chapter 10.4, pp.611--623.
\chapter{Rational Choice Theory in the Social Sciences}
\label{sec:p17-rationality-10-4}

Rational choice assumptions can be ingredients of an explanation without
being the whole explanation. A social scientist who wants to explain whether
a group supplies a collective good needs not only a theory of individual
choice but also an account of who can act, what they know, how their actions
affect one another and how those actions produce the group outcome.
Raub's chapter develops this model-building use of rational choice, chiefly
in sociology rather than economics
\parencite[pp.~611--623]{Raub2021RationalChoiceSocialSciences}.
The explanatory unit is a conditional social-science model, not a preference
aggregation rule and not a universal claim that people calculate selfish
monetary returns.
\index{rational choice theory!social explanation}
\index{micro--macro explanation}\index{methodological individualism}

Raub uses \emph{RCT} for social-science theories that employ rational choice
assumptions together with other assumptions. The abbreviation therefore does
not here mean only the axioms of rational choice. His historical orientation
links applications in political science to Arrow, Buchanan and Tullock,
Downs, Olson and Ostrom; in sociology it links Coleman to earlier behavioral
sociology and social exchange, and to European structural individualism or
explanatory sociology. These are the source's disciplinary connections, not
a claim that those approaches have identical premises or that their authors
all originated one theory. Raub explicitly leaves a fuller history of ideas
outside his selected treatment
\parencite[pp.~611, 620, note~1]{Raub2021RationalChoiceSocialSciences}.

\section{The explanatory unit and its assumptions}
\label{sec:p17-rat104-explanatory-unit}

\begin{hbdefinition}{def:p17-rat104-rct}{Action assumptions within a social-science model}
An RCT explanation relates a specified social situation to actors' choices
and those choices to social outcomes. Its \emph{micro-level} describes actors,
including their feasible actions, preferences, beliefs and behavior. Its
\emph{macro-level} describes the social system of which those actors are
parts, such as a family, organization, city, school or market.
Methodological individualism is the model-building requirement to account
for system-level phenomena through actors and their relations; it is not
the assertion that social conditions do not matter.

The levels are relative to the question. Employees can be actors within a
firm; firms can instead be corporate actors within an interorganizational
network. Treating a firm as a unitary actor is a substantive simplifying
assumption, not a conclusion from its legal identity. Small dyads or triads
are systems too: macro does not mean numerically large.
\end{hbdefinition}

These distinctions preserve Raub's broad use of social system and his
corporate-actor examples
\parencite[pp.~611--612, notes~1--2 on p.~620]{Raub2021RationalChoiceSocialSciences}.
Corporate action as an explanatory simplification is distinct from asking
whether a group possesses a genuinely collective standpoint; compare
\cref{def:p17-rat102-two-standards}.

\begin{hbdefinition}{def:p17-rat104-assumption-types}{Rationality, motives, information and social structure}
Four assumption types must be distinguished.
\begin{enumerate}
\item \emph{Rationality assumptions} constrain the relation between
      preferences, beliefs and choices: for example, preference transitivity,
      expected-utility choice on an appropriate lottery domain, or a stated
      game-theoretic equilibrium condition.
\item \emph{Substantive motivational assumptions} say what consequences the
      actor values, how much, and with what variation across actors. Caring
      about others is compatible with a coherent utility representation.
\item \emph{Information assumptions} specify what actors observe, know or
      believe, including uncertainty about consequences or others' motives.
\item \emph{Structural assumptions} specify opportunities, constraints,
      interaction, institutions and the mechanism producing social outcomes.
\end{enumerate}
An empirical explanation requires all relevant types. Transitivity alone
does not tell us whether an actor prefers one concrete alternative to
another, and utility maximization alone does not identify the utilities.
\end{hbdefinition}

Raub's standard micro-model combines rational choice with self-regarding
preferences and often perfect or complete information. These are distinct
premises, not one indivisible meaning of rationality. Rational choice need
not mean that utility equals own money, that information is complete, or
that behavior is deterministic. Mixtures in the worked example below are
explicitly probabilistic. Nor does an as-if utility representation imply
conscious arithmetic in the actor's mind
\parencite[pp.~616--617, 619]{Raub2021RationalChoiceSocialSciences}.
The finite lottery representation in \cref{thm:p17-finite-vnm} specifies its
own axioms; \cref{def:p17-vnm-interpretation} separates representation from
empirical conformity and interpersonal welfare.

\section{Connecting conditions, actions and outcomes}
\label{sec:p17-rat104-links}

Coleman's diagram separates four kinds of propositions and four relations.
Its usefulness lies in exposing missing explanatory links, not in making
them true simply by drawing arrows
\parencite[pp.~611--614, figs.~10.4.1--2]{Raub2021RationalChoiceSocialSciences}.

\begin{hbdefinition}{def:p17-rat104-coleman}{Bridge assumptions and transformation rules}
Write \(A\) for macro-conditions, \(B\) for micro-conditions,
\(C\) for micro-outcomes and \(D\) for macro-outcomes. Here these letters
name classes of propositions, not events in a probability space.
\begin{center}
\begin{tabular}{p{0.20\textwidth}p{0.72\textwidth}}
\(A\longrightarrow B\) & Bridge assumptions: how social conditions
change feasible alternatives, incentives or information.\\
\(B\longrightarrow C\) & Action theory: how actors choose under those
conditions and the specified behavioral assumptions.\\
\(C\longrightarrow D\) & Transformation rules: how individual actions
combine to generate a system outcome.\\
\(A\longrightarrow D\) & The macro relation to be explained or predicted,
often initially identified as an empirical association.
\end{tabular}
\end{center}
Explaining \(C\) needs the initial macro-conditions, bridge assumptions and
action theory. Explaining \(D\), or a regularity linking \(A\) and \(D\),
additionally needs transformation rules. A macro association by itself does
not establish this intervening mechanism or its causal direction.
\end{hbdefinition}

For instance, an electoral institution can change incentives to vote or
choose a party, and can also change the translation of unchanged votes into
seats. The first is a bridge effect; the second is a transformation effect.
Conflating them can make an institutional outcome look like a change in
voter preferences when none is assumed. This is Raub's qualitative example,
not an estimated effect of a named electoral reform
\parencite[p.~618]{Raub2021RationalChoiceSocialSciences}.

\begin{hbmethod}{def:p17-rat104-model-building}{Constructing and testing a micro--macro explanation}
\begin{enumerate}
\item Specify the explanandum: an individual action, a system outcome, or a
      relation between system conditions and outcomes. Fix the system,
      actors, times and comparison domain.
\item State the structural conditions and bridge assumptions. Give the
      feasible choices, outcome consequences and information available when
      choices are made; a payoff table alone need not specify information.
\item State the action theory and substantive preference or belief premises
      separately. If several equilibria satisfy the theory, disclose the
      additional selection assumption used for a prediction.
\item Supply the transformation rule. Distinguish an individual's action
      probability from the system's probability of success.
\item Derive both levels' implications, retaining the hypotheses used.
      Compare a nearby model that relaxes a consequential assumption.
\item Test the relevant implications and, where possible, the assumptions
      or links. A mismatch can concern motives, information, structure,
      behavior, transformation or selection; it does not identify its own
      cause merely by occurring.
\item Refine problematic assumptions stepwise and derive new implications.
      Assess whether the refinement transfers across situations or merely
      fits the original observation after the fact.
\end{enumerate}
This procedure reconstructs Raub's model-building argument; it is not an
algorithm that guarantees an adequate social explanation.
\end{hbmethod}

Formalization matters because opposing qualitative effects need not reveal
their net effect. Tractability and realism can pull in different directions:
simple assumptions permit deductions but may omit a material mechanism.
The method of decreasing abstraction begins with explicit simplifications
and adds complexity when conclusions depend sensitively on them or fare
poorly against evidence. Raub presents Coleman's advice to invest especially
in bridge assumptions and transformation rules, while retaining parsimonious
micro-assumptions where possible. This is a research strategy, not proof that
micro-errors always cancel. Reusing a common action-theoretic core across
middle-range models supports coherence without turning every model into one
unconditional macro-law
\parencite[pp.~614--615, 619--620]{Raub2021RationalChoiceSocialSciences}.

\section{The Volunteer's Dilemma: individual choice and collective provision}
\label{sec:p17-rat104-vod}

Raub's main example is a collective good that cannot be withheld from
noncontributors. A costly contribution has an individual incentive problem
when an actor expects someone else to provide the benefit. Nonexcludability
alone, however, does not imply a universal negative group-size law: selective
incentives, the production function and other conditions matter. The
Volunteer's Dilemma isolates one exact configuration
\parencite[pp.~612--615]{Raub2021RationalChoiceSocialSciences}.

\begin{hbdefinition}{def:p17-rat104-vod}{The source's symmetric provision game}
There are \(N\geq2\) actors and common cardinal benefit and cost
\(U>K>0\). Each chooses \(\mathrm{CONTR}\) or \(\text{DON'T}\).
The good is supplied if and only if at least one actor contributes.
Every contributor pays the full cost \(K\), even if others contribute;
extra contributions do not increase the benefit. Payoffs are
\[
\begin{array}{c|cc}
 &\text{no other contributor}&\text{at least one other contributor}\\ \hline
\mathrm{CONTR}&U-K&U-K\\
\text{DON'T}&0&U
\end{array}
\]
Choices are simultaneous: when choosing, an actor does not observe any
other actor's choice. Concretely, one may order the decisions in a game
tree while giving each actor one information set containing all preceding
histories; no preceding move is revealed. Payoffs and this structure are
shared and known in the model. Use independent randomization, with
\(p_i\in[0,1]\) the probability that actor \(i\) contributes.
Then
\[
 P(p_1,\ldots,p_N)=1-\prod_{i=1}^N(1-p_i),\qquad
 E_i(p_1,\ldots,p_N)=U P(p_1,\ldots,p_N)-Kp_i.
\]
Here \(P\) is provision probability and \(E_i\) expected utility,
not a moral evaluation of helping. The cost and benefit are consequence
utility levels in the specified model; no interpersonal utility sum is
needed for Pareto comparisons.
\index{volunteer's dilemma}\index{collective good!provision probability}
\end{hbdefinition}

The explicit information-set description fills the boundary identified in
Raub's note~4: a normal form alone does not fully specify the information
structure. Independent randomization is a separate probability assumption,
not the preference-independence axiom in expected-utility theory. The
bystander analogy concerns witnessing an accident or crime, relief from
someone helping and the cost of one's own intervention. It motivates the
model; it is not evidence that every emergency has identical costs, knowledge,
opportunities or equilibrium behavior
\parencite[pp.~613--614, note~4 on p.~620]{Raub2021RationalChoiceSocialSciences}.

\begin{hbtheorem}{thm:p17-rat104-vod-equilibria}{proved=annex:p17.rat104-vod-equilibria}{Equilibria and the conditional group-size effect}
In \cref{def:p17-rat104-vod}, there are exactly \(N\) pure Nash
equilibria, each with exactly one contributor. Each supplies the good
with certainty. There is exactly one symmetric independent mixed Nash
equilibrium, at which every actor contributes with probability
\[
 p^*=1-\left(\frac KU\right)^{1/(N-1)},\qquad
 P^*=1-(1-p^*)^N
     =1-\left(\frac KU\right)^{N/(N-1)}.
\]
Every actor's equilibrium expected utility is \(U-K\). Holding \(K/U\)
fixed, both \(p^*\) and \(P^*\) strictly decrease as integer \(N\)
increases; their limits are respectively \(0\) and \(1-K/U\).
This is a prediction conditional on selecting the symmetric equilibrium,
not a property of every equilibrium or an observed population frequency.
\end{hbtheorem}

Symmetry of the game does not eliminate asymmetric equilibria. Each actor
prefers another actor to be the sole volunteer: the nonvolunteer obtains
\(U\), the volunteer \(U-K\). Raub treats the symmetric mixture as a
plausible selection in the absence of an agreed assignment, not as a result
that rationality uniquely selects it. At a fixed positive individual
probability, adding potential volunteers raises provision; equilibrium
adjustment instead lowers each individual's probability enough that provision
falls. The transformation rule is what determines which effect wins
\parencite[pp.~613--614]{Raub2021RationalChoiceSocialSciences}.

\begin{hbexample}{def:p17-rat104-numeric}{The same incentives give different outcomes under different selection}
Take \(U=10\) and \(K=5\). For \(N=2\), the symmetric equilibrium has
\(p^*=1/2\) and \(P^*=3/4\). For \(N=3\),
\(p^*=1-1/\sqrt2\) and \(P^*=1-1/(2\sqrt2)\), approximately
\(0.2929\) and \(0.6464\). Each actor's expected utility is 5 in both
cases. The two-actor sole-volunteer equilibria instead give utilities
\((5,10)\) or \((10,5)\) and certain provision.
Thus increasing group size does not itself imply falling provision without
the selection and structural assumptions. These numbers are calculations
in the supplied game, not experimental estimates.
\end{hbexample}

\subsection{Joint improvement and the domain of efficiency}

Raub compares the equilibrium with a more cooperative symmetric mixture,
retaining the same game and independent draws
\parencite[p.~615]{Raub2021RationalChoiceSocialSciences}.

\begin{hbtheorem}{thm:p17-rat104-cooperative-mixture}{proved=annex:p17.rat104-cooperative-mixture}{The cooperative symmetric profile and its precise Pareto boundary}
Set
\[
 p^{**}=1-\left(\frac{K}{NU}\right)^{1/(N-1)},\qquad
 a=1-p^{**},\qquad q=1-\frac{N-1}{N}a.
\]
Then \(0<p^*<p^{**}<1\). The symmetric profile with contribution
probability \(p^{**}\) uniquely maximizes each actor's common expected
utility among symmetric independent profiles, with value
\[
 E_i=U-qK>U-K,\qquad 0<q<1.
\]
It is Pareto efficient among all independent mixed profiles, including
asymmetric ones: none makes everyone at least as well off and somebody
strictly better off. It is not a Nash equilibrium. Against opponents using
\(p^{**}\), an actor strictly prefers not to contribute.
\end{hbtheorem}

Here defection and cooperation can label the mixtures \(p^*\) and
\(p^{**}\) respectively. They do not identify \(\text{DON'T}\) and
\(\mathrm{CONTR}\) as the dominant actions of a prisoner's dilemma.
In this game all abstaining is not an equilibrium: a sole volunteer gains
\(U-K>0\). All contributing is not an equilibrium either. The dilemma is
that moving together to a more cooperative distribution improves everyone,
while each actor has a unilateral incentive to avoid the cost once others
make that move. This preserves the source's mixed-strategy comparison
without treating all social dilemmas as one payoff structure.

Pareto efficiency depends on the allowed profiles. If a coordinating device
uniformly assigns exactly one actor to volunteer, each actor obtains
\(U-K/N\). The assignment guarantees provision with one total contribution;
it gives each actor \(K(N-1)p^{**}/N>0\) more expected utility than the
independent \(p^{**}\) profile. This correlated assignment lies outside
the theorem's independent-mixture domain. Its desirability also does not
show that actors will comply with it: a recommended volunteer needs reasons
or incentives to accept the recommendation, and an implemented institution
needs its own model.

\begin{hbtheorem}{thm:p17-rat104-provision-correction}{proved=annex:p17.rat104-provision-correction}{Provision under the cooperative mixture need not increase with group size}
Holding \(r=K/U\in(0,1)\) fixed, provision under \(p^{**}\) is
\[
 P^{**}(N)=1-\left(\frac rN\right)^{N/(N-1)}.
\]
It is not increasing for every allowed parameter and integer group size.
For \(r=1/2\),
\[
 P^{**}(2)=\frac{15}{16}
   >1-\left(\frac16\right)^{3/2}=P^{**}(3).
\]
For its continuous extension to real \(N>1\), the derivative is positive
exactly when \(\log(N/r)<N-1\), negative when the inequality reverses,
and zero at equality. Consequently it eventually increases and converges
to 1 as group size grows. For \(r\geq2/e\) it is strictly increasing
on integer \(N\geq2\).
\end{hbtheorem}

The counterexample is within the declared \(U>K>0\) domain;
the increase criterion is therefore essential. The equilibrium result about declining \(P^*\)
remains unchanged. Neither model calculation settles actual group-size
effects in emergencies or collective-action settings
\parencite[pp.~614--615, note~5 on p.~620]{Raub2021RationalChoiceSocialSciences}.

\section{Variants change more than one dimension}
\label{sec:p17-rat104-variants}

Raub distinguishes variation in the micro-model from variation in the
social structure and micro--macro links. A criticism of one combination
does not refute every combination
\parencite[pp.~616--618]{Raub2021RationalChoiceSocialSciences}.

On the micro side, prospect theory, bounded-rationality accounts,
dual-process or two-selves accounts, and backward-looking learning modify
the action theory. Esser's frame-selection model and Lindenberg's social
rationality are sociological approaches informed by such alternatives.
Coleman's diagram itself is neutral about the behavioral theory placed
between \(B\) and \(C\); it does not prove that these alternatives are
equivalent or empirically adequate. Related treatments here include
\cref{sec:p17-rationality-8-3,sec:p17-rationality-8-5}.

Other variants retain rationality while adding other-regarding preferences,
heterogeneity in their strength and incomplete information about others'
utilities. This changes motives and information, not necessarily the
coherence standard. Raub reports inequity-aversion and related models that
address dilemmas, ultimatum and dictator interactions and markets, including
the possibility that the same motives generate cooperative behavior in one
structure and seemingly selfish behavior in another. No single fitted
preference parameter or effect size is established by that survey statement
\parencite[pp.~616--617]{Raub2021RationalChoiceSocialSciences}.

The objection is underdetermination: if one chooses a new utility function
after seeing each act, almost any behavior can be made to look optimal.
Raub's safeguards are parsimony, using the same assumptions across a range
of situations, and deriving new testable predictions rather than fitting
only familiar results. The Stigler--Becker strategy instead holds fundamental
preferences relatively stable and explains changes through opportunities
and constraints; Raub relates social-production-function theory to that
strategy. These are rival modeling commitments, not evidence that tastes
are always invariant or that concern for others is unscientific
\parencite[pp.~616--617]{Raub2021RationalChoiceSocialSciences}.

On the macro side, Raub treats the perfect-market model as a benchmark for
deriving system outcomes from individual exchange. The cited existence and
welfare results depend on a fully specified market's hypotheses; the
unspecified phrase perfect market does not license an equilibrium-existence
or Pareto-optimality theorem here. Their role in his argument is to show
why deriving macro consequences is substantive rather than obvious. Outside
that benchmark, strategic interdependence may matter directly.
Granovetter's embeddedness concerns continuing relations and networks;
Raub reads his emphasis on social structure as compatible with keeping
rationality while refining constraints, information and interaction
\parencite[pp.~617--618]{Raub2021RationalChoiceSocialSciences}.

Networks, institutions and organizations may be \emph{exogenous} conditions
affecting choices or their combination. They may instead be \emph{endogenous}
outcomes of the actions being explained. Formal institutions include legal
rules; informal ones include conventions and behavioral codes. An endogenous
account can connect one period's \(D\) to the next period's \(A\), but
the transitions need explicit rules. A static payoff matrix is not a theory
of network formation, institutional emergence or evolutionary adjustment
\parencite[p.~618]{Raub2021RationalChoiceSocialSciences}.

\section{Interdependence, unintended consequences and empirical limits}
\label{sec:p17-rat104-evidence}

When an outcome depends on other actors as well as one's own action, intended
ends need not match realized system consequences. Rational choice does not
remove this possibility. Raub contrasts beneficial unintended consequences,
illustrated by his perfect-market benchmark, with unfavorable ones in social
dilemmas. In the Volunteer's Dilemma, maximizing conditional individual
payoffs can yield the independently Pareto-inferior symmetric equilibrium
without anyone intending the inferior provision pattern. A value judgment
about a realized consequence is distinct from a model of its production
\parencite[pp.~618--619]{Raub2021RationalChoiceSocialSciences}.

Raub presents game theory as a tool for specifying interdependent choice,
including extensive, repeated and incomplete-information games rather than
only complete-information normal forms. His connection to Weber's social
action emphasizes orientation toward others' behavior; it does not say
that every socially meaningful action satisfies one equilibrium concept.
Learning and evolutionary games can model adjustment without assuming
deliberation of the same kind; compare
\cref{sec:p17-rat93-evolutionary-games}.

\begin{proofboundary}{def:p17-rat104-evidence-boundary}{A survey report is not an original-study estimate}
Raub reports experimental work in which individual volunteering declines
with group size but collective provision does not show the predicted decline,
and identifies Tuti\'c's 2014 procedurally rational model as an alternative.
The chapter supplies no study-specific task version, participant population,
assignment design, sample size, measured effect or uncertainty estimate for
that report. It is retained here as a source-qualified limitation on the
symmetric-equilibrium prediction, not as a new universal empirical law or
as an independently verified numerical result. The alternative's complete
solution concept is not reconstructed from this brief mention
\parencite[p.~615]{Raub2021RationalChoiceSocialSciences}.

Likewise, Raub's reports of laboratory and field anomalies, broad empirical
applications and market deviations summarize cited literatures. They do
not supply a common population or one measured effect. An empirical claim
about a specified population requires the evidence record of
\cref{def:p19-evidence-record}, including task, design, measures,
uncertainty and generalization limits. None of these reports by itself
establishes what actors morally ought to do
\parencite[pp.~616, 619--620]{Raub2021RationalChoiceSocialSciences}.
\end{proofboundary}

The source's discussion connects theory and research in both directions.
Empirical regularities identify explananda, not just tests of a prior model.
Laboratory and field experiments, surveys, observational designs and
vignettes can test related hypotheses under complementary conditions.
Agreement may support robustness but does not erase differences in task,
population or identification. Raub also mentions conditional-logit and
network-dynamics work as ways to integrate action assumptions into a
statistical model; naming those approaches does not establish fit or
identification for an unspecified data set
\parencite[pp.~619--620]{Raub2021RationalChoiceSocialSciences}.

Empirical evaluation can target predictions alone or also the preferences,
beliefs and links used to derive them. A rational model may serve as a
benchmark: systematic deviations then become explananda for a refined
model. But neither the argument that prediction errors cancel in aggregate
nor the argument that competitive institutions generically eliminate
nonrational behavior is assured. Raub records both arguments and the
objection that deviations can be patterned and can survive aggregation.
He also records disagreement over comparative-static directional predictions
versus quantitative point predictions. A model that gets an individual
direction right can get the collective effect wrong, as the distinct
\(p\) and \(P\) in the worked example make explicit
\parencite[pp.~615, 619--620]{Raub2021RationalChoiceSocialSciences}.

\section{Subject and task lookup boundaries}
\label{sec:p17-rat104-lookup}

\begin{hbmethod}{def:p17-rat104-lookup}{Choosing the explanatory question, rather than a neighboring formal task}
For explaining how a social condition produces an individual and collective
outcome, use the micro--macro model-building procedure above. A successful
case specifies bridge assumptions, action theory, transformation and
selection, and derives the conditional Volunteer's Dilemma prediction.
A failed case infers provision from an individual's contribution probability
alone, assumes symmetry selects an equilibrium without further premise,
or transfers the calculated decline to every real group.

For evaluating an individual lottery menu, use
\cref{def:p17-expected-utility-choice}. Its utility representation does not
supply social opportunities or an aggregate mechanism. For testing unilateral
strategic deviations, use \cref{def:p17-nash-equilibrium}; equilibrium
does not imply joint efficiency, intention or observed play. For forming a
social ranking from individual rankings, use
\cref{sec:p17-social-choice}; its aggregation axioms do not explain how a
network or incentive affects behavior. For consistency of group judgments,
use \cref{sec:p17-judgment-aggregation}; a consistent judgment set need
not provide a collective good. For cooperation, coordination and group
agency as standards of collective rationality, use
\cref{sec:p17-rationality-10-2}; the present explanatory model does not
settle a group's normative authority or ontology.
\end{hbmethod}

Raub's further-reading map places behavioral and experimental game theory,
learning and evolutionary games, analytical sociology and agent-based
modeling alongside RCT because they can connect incentive-driven or
goal-directed action to social processes. Their overlap in question does
not collapse their assumptions. The contribution of this treatment is the
explicit explanation unit, its complete worked formal example and its
conditional evidence boundaries, not a claim that one rational-choice
model exhausts social science
\parencite[p.~620]{Raub2021RationalChoiceSocialSciences}.
